\clearpage
\chapter*{第四部分结语}
\addcontentsline{toc}{chapter}{第四部分结语}
\label{part:04-closing}

\begin{center}
  \includegraphics[width=\linewidth,height=.35\textheight,keepaspectratio]{assets/part-art/part-4-closing.png}
  \par\vspace{4pt}
  {\scriptsize\color{Muted}取意于卡米耶·毕沙罗《蒙马特大道，春天的早晨》\par}
  \BookEpigraph{part4}
\end{center}

第四部分把购买服务展开为一套共同运行的系统。入口根据目录找到购买权威，库存组保存资源预留，支付平台掌管实际收款，消息域传递事件，消费事务库保存履约与发运意图。每个组件有自己的接口、保存状态和故障处理，原购买的关联把这些位置接到同一条业务链上。沿链查看一次处理，可以找到当前决定由谁作出、结果保存在哪里、接下来的工作由谁承担。

购买权威保存完整身份到实际 P/C/T 的绑定，以及随后形成的原结果、订单与 EV。分区复制状态机按提交顺序共同修改这些关系，也可以选择由分布式 SQL 的事务与约束承担相应修改。复制保存确认过的决定，快照包含会话、结果、编号进度和未决责任。读取原结果时，服务沿原绑定返回当次事实，当前履约与退款查询则继续读取各自的新进度。

控制面保存逻辑分区当前的位置和执行资格。目录更新以后，入口缓存逐步刷新，实际接收写入的组核验归属与代次。协调服务的版本、会话和通知帮助应用接续配置，业务资源再在实际修改处检查资格。查询缓存则从权威事实形成可重建观察，过期、淘汰和节点切换以后，按保存的来源与版本重建所需状态。

外部付款从受理时保存的 C/T 继续。工作者根据领取代次安排发送、查询与回填，发送使用已保存的平台、商户范围、环境、动作和 \code{pay-C}。X1 已经完成而本地尚未取得确定证据时，付款保持未知。查询、可信回调和对账逐步核对原收款，权威随后固定结果、订单与事件。平台只提供时点查询，或原键保护期和在途请求尚有歧义时，原 C 留在持续核查路径，调度依据已保存状态安排后续查询。

联合库存预留在各库存组之间组织共同决定。参与者保存准备占用，协调者持久保存 commit 或 abort，恢复者把原决定送到各组。参考流程等待全部必需参与者完成以后，才开启原 C/T 的付款资格。库存占用、协调者决定和付款资格都进入各自的复制状态，接班者据此继续传播决定或推进付款。

付款形成的原 EV 可以经过重复发布和重投。RabbitMQ 的路由、队列与确认，Kafka 的分区日志、事务与消费位置，分别规定消息域内部的处理。消费者在本地事务中保存 I/B/F、结果和 dispatch 意图。消息保留期有限时，持久消费交接回执把 F 的接受事实送回购买权威，更新履约进度并结束原事件的补发责任，实际发运则沿原 dispatch 继续查询外部平台，把确定结果接回相应的履约记录。

退款在这条链上建立新的 RF，关联原 C 与 X1。收款 3000 分，成功退款 1000 分，另有 1500 分未知，后续最多再占用 500 分。调度转入人工核查以后，未知的 1500 分继续保留占额，取得确定证据后再转换或释放。财务、客服与用户查询由此可以同时解释原购买成功、当前净资金与退款进度。

存储与计算组件也参与这些保存和接续。块与对象保存图像、日志和文件，发布协议确定哪个版本或文件集合已经生效。批处理以固定输入组织重执行，用提交结果确定有效输出；任务、Actor 与函数运行时则按各自规则调度和恢复执行。持久工作流保存步骤位置，付款步骤从原平台关联继续。推测执行、共识优化和地域读协议调整具体等待，部署为相应日志、检查点、资格和证据安排保存与恢复资源。

完整分区迁移把这些业务关系交给新的承载组。参考首版先闭合 prepared 库存与未完成协调者事务，再冻结图像并切换目录，已受理的外部责任随原关联交接。旧备份恢复则先找回旧身份的绑定和较新的外部证据，或者持久关闭歧义范围，再开放相应入口。缓存重建、快照下载、日志接续和任务扫描分别占用网络、磁盘与执行能力，恢复放量沿这些依赖分批推进。

容量设计最后把整套架构连到实际运行。平台恢复使 T 下降，同时可能产生更多 EV 与消费积压。调度按整条链的净排空能力安排新流，为原结果查询、确定证据回填、事务决定恢复和事件交接保留预算。观测从原身份找到各段责任，再显示积压年龄、最近证据和下一步动作，值守者便能据此限制新增负担、调整恢复份额或发起核查。
\FloatBarrier
